\documentclass[11pt]{article}

\usepackage[margin=0.50in]{geometry}
\usepackage[T1]{fontenc}
\usepackage{lmodern}
\usepackage{microtype}
\usepackage{enumitem}

\setlength{\parindent}{0pt}
\setlength{\parskip}{0.45em}
\setlist[itemize]{leftmargin=1.45em,itemsep=0.25em,topsep=0.2em}
\setlist[enumerate]{leftmargin=1.65em,itemsep=0.55em,topsep=0.25em}

\newcommand{\location}[1]{%
  \par\nobreak\smallskip
  \textit{Locations in the revised manuscript:} #1.
}

\begin{document}

\begin{center}
  {\Large\bfseries Summary of Changes}\par
  \vspace{0.25em}
  {\large ASE 2026 Major Revision}\par
  \vspace{0.15em}
  \textit{DesignFix: Context-Augmented Repair of Design Rule Violations}
\end{center}

\section*{1. Mandatory changes requested in the metareview}

\begin{enumerate}

\item \textbf{We added a lightweight autonomous-agent baseline on the
directly matched DesignFix-condition tasks from the human study and used it to
identify DesignFix's distinct contribution.}

We added RQ4 and a new evaluation using Cline, a general-purpose autonomous
in-IDE coding agent. 

\location{Introduction (RQ4 and contributions); Related Work, ``LLM Agents in
Software Development''; Evaluation III, ``Lightweight Agentic Baseline'';
Discussion, ``Autonomy solved localization, not design intent'' and
``DesignFix and autonomous agents are complementary''; Conclusions}

\item \textbf{We substantially revised the threats to validity to make the
user-study design and the scope of its results explicit.}

The revised manuscript now states that the eight user-study rules were selected
purposively rather than randomly. We explicitly acknowledge that
conditioning task selection on observed outcomes and author judgment may bias
the estimated condition effect.

The revised threats section additionally discusses possible repository-learning
carryover, self-selected browser-LLM use, prompt-development overlap, the
single-rater assessment, the limited participant count, and the descriptive
nature of the agent evaluation.

\location{Evaluation II, ``Methodology---Tasks and task selection''; Threats to
Validity, particularly Construct Validity, Internal Validity, and External
Validity}

\item \textbf{We added an explicit discussion of how modern autonomous in-IDE
agents affect the relevance and future role of DesignFix.}

The Introduction now motivates the work in a setting where developers may use
both conversational assistants and autonomous agents such as Cline and Claude
Code.  A new related-work subsection reviews LLM agents in software
development. 

The Discussion now distinguishes capabilities likely to improve with stronger
agents from
inputs that autonomy does not create by itself.  This makes the future-proofing claim
specific: as agents improve, DesignFix's checker, context extraction, example
retrieval, and repair validation can be exposed as agent-callable services
rather than replaced by the agent.

\location{Introduction; Related Work, ``LLM Agents in Software Development'';
Evaluation III; Discussion; Conclusions}

\end{enumerate}

\section*{2. Clarifications promised in the rebuttal}

\begin{enumerate}

\item \textbf{Baselines and information available in each condition.}
Evaluation~I now defines its Baseline as direct non-agentic LLM use: the model
receives only the rule description and violating snippet, approximating a
developer pasting a checker warning and flagged code into a conversational
assistant. 

Evaluation~II is now consistently described as a comparison \emph{with versus
without DesignFix}, not as LLM assistance versus no LLM assistance. 

\location{Evaluation I, ``Methodology---Prompt Development'' and Results;
Evaluation II, ``Methodology---Study Design and Procedure''; Discussion,
``The measured gains are over general LLM assistance, not over no assistance''}

\item \textbf{Problem prevalence and practical significance.}
We added empirical motivation from prior works.

\location{Introduction; Background, ``Design rule violations are common and
consequential''}

\item \textbf{Rule construction, satisfying examples, data provenance, and
checker limitations.}
The Background now gives an end-to-end example connecting a natural-language
DTO rule, its XPath quantifier and constraint over srcML, an actual violation,
and a satisfying example.  It precisely defines a satisfying example,
and explains the natural-language intent.  We also explain why the complete file and
the separately labelled.

We corrected the dataset description from two codebases to two corpora.
CrowdCode is the studied application, and its rules and seeded violations were
inherited from prior work.  Android Architectural Rules is a purpose-built
evaluation corpus; its violations were introduced into those fixtures by
the authors. 

\location{Background, particularly ``Structure of a design rule'' and ``A
concrete example, end to end''; DesignFix System Design, ``System
Architecture'' and ``Context Extraction''; Evaluation I,
``Methodology---Dataset''}

\item \textbf{Evaluation rubrics and assessment procedure.}
We now define Code Fix Success rate, Surrounding Code Pass Rate and Explanation Correctness rigorously and strictly. Our current definition of surrounding code pass rate is correct and differs from the definition mentioned in the rebuttal.

\location{Evaluation I, ``Methodology---Metrics''; Evaluation II,
``Methodology---Measures''; Threats to Validity, Construct Validity}

\item \textbf{Novelty and scope of generalization.}
The contribution claims have been revised to distinguish the use of familiar
components from the empirical and workflow contributions.  The paper now
centers its contribution on determining which checker-derived and
project-specific evidence enables reliable design-rule repair, integrating
that evidence into a reviewable workflow, and evaluating the complete workflow
against both browser-LLM-assisted developers and an autonomous agent. We also narrowed down RQ1.


\location{Introduction, contributions; DesignFix System Design; Threats to
Validity, External Validity; Discussion}

\end{enumerate}

\section*{3. Corrections made during final data verification}

During revision, we rechecked the underlying study records and corrected two
values stated in the rebuttal.  The mean time before the first code edit is
5.8 minutes without DesignFix and 2.3 minutes with DesignFix, replacing the
preliminary 7.2- and 1.8-minute values.  Across all six matched
model--temperature conditions, the gain from adding a same-project satisfying
example is 7--10 percentage points, replacing the preliminary 8--13-point
range; at temperature 0.5, the gains are 8 points for GPT-4o-nov and 7 points
for GPT-4o-aug.  The revised manuscript uses the reverified values
consistently.

\location{Evaluation I, Methodology and Key Findings; Evaluation II,
Qualitative Findings; Threats to Validity, External Validity}

\section*{4. Reproducibility and presentation}

The manuscript now uses the ACM review format, includes an enlarged system
diagram and a referenced end-to-end rule listing, and cites the anonymized
DesignFix replication package in the Data Availability Statement.

\location{Background; DesignFix System Design; Data Availability Statement}

\end{document}
